
\needspace{4\baselineskip}
\color{uniblue}{\subsection{Контроль наличия напряжения до ввода (КННвв)}}
\color{black}

\begin{enumerate}[label=\arabic{section}.\arabic{subsection}.\arabic*, labelsep=4pt, leftmargin=0pt, itemindent=57pt]

\item
КННвв контролирует наличие напряжения до ввода с дополнительным контролем симметрии по $U_2$. В устройстве контроль наличия напряжения до ввода представлен функцией «КННвв». Функция используется для АВР, ВНР, ЗМН и АПВ с контролем напряжения на вводе, а также для оперативного включения.

\item
Пусковые условия функции формируются без выдержки времени при превышении минимального или максимального из подведенных линейных напряжений выше уставки в зависимости от того сколько фаз будет подведено к устройству.

\item
Функция предусматривает два режима работы, определяемые уставкой SGF2 «Реж\_пуска»:
\begin{itemize}
    \item измеренный;
    \item внешний.
\end{itemize}

\item
Алгоритм работы функции, когда уставка SGF2 «Реж\_пуска» выставлена в положение «Измеренный»:
\begin{itemize}
    \item если к устройству подводятся три фазных или два линейных напряжения, то уставка SGF3 «Контроль\_ТН» должна быть выставлена в положение «3 фазы». В этом случае пусковые условия будут формироваться, когда все три линейных напряжения будут превышать значение уставки с отсутствием срабатывания ИО ОП выше уставки ;
    \item если к устройству подводится только одно линейное напряжение, то уставку SGF3 «Контроль\_ТН» необходимо выставить в положение «2 фазы». В этом случае пусковые условия будут формироваться при превышении максимального из линейных напряжений выше уставки без учета срабатывания ИО ОП.
\end{itemize}

\item
Предусмотрена блокировка функции при срабатывании ИО обратной последовательности напряжения выше уставки.

\item
Если уставка SGF2 выставлена в положение «Внешний», работа функции осуществляется при наличии активного внешнего дискретного сигнала «Наличие Uвв» без учета срабатывания измерительных органов.

\item
Ввод функции в работу на этапе параметрирования устройства осуществляется программным переключателем SGF1 «Ввод\_функции». Информация о введенной функции отображается выходным сигналом «КННвв: Ввод».

\item
Функция может быть выведена из работы оперативно путем активации сигнала «ОВ КННвв», что характеризуется активным состоянием сигнала «КННвв: Оперативный вывод» на выходе функции.

\item
Алгоритм работы функции выполнен в соответствии с рисунком \ref{knnvv:img1}. В таблице \ref{knnvv:tbl1} приведены параметры, необходимые для настройки функции.

\vspace{3mm}
\begin{figure}[H]
\centering
\includegraphics[width=0.9\textwidth,height=0.9\textheight,keepaspectratio]{\fbpath/306. КННвв/_latex/img/knnvv.pdf}
\captionof{figure}{Функциональная схема логики КННвв}\label{knnvv:img1}
\end{figure}

\small
\begin{longtable}{|>{\centering\arraybackslash}m{5.3cm}|>{\centering\arraybackslash}m{3.3cm}|>{\centering\arraybackslash}m{4.2cm}|>{\centering\arraybackslash}m{1.8cm}|>{\centering\arraybackslash}m{1cm}|}
\caption{Параметры для настройки КННвв\hfill\vspace{-0.5\baselineskip}}\label{knnvv:tbl1}\\ 
\hline
\rowcolor{gray!30}
Параметр (Параметр на ИЧМ) & Условное обозначение на схеме & Значение/ Диапазон & Единица измерения & Шаг \\ 
\hline
\endfirsthead
\caption*{\hspace{3pt}\emph{Продолжение таблицы \ref{knnvv:tbl1}\hfill\vspace{-0.5\baselineskip}}} \\ % сделано по ГОСТ 2.105 п.6.8.7
\hline
\rowcolor{gray!30}
Параметр (Параметр на ИЧМ) & Условное обозначение на схеме & Значение/ Диапазон & Единица измерения & Шаг \\ 
%\hline
\endhead
%\hline
\endfoot
\hline
\endlastfoot
%==+t1*PTOV1|VV_VOLPRES2
\centering Ввод функции в работу (Ввод\_функции)  & \centering SGF1 & \centering Не предусмотрено \\ Предусмотрено & \centering -- & \centering \arraybackslash -- \\
\hline
\centering Напряжение срабатывания (Uср)  & \centering U> & \centering 50,0...150,0 & \centering В & \centering \arraybackslash 0,1 \\
\hline
\centering Напряжение срабатывания ОП (U2ср)  & \centering U2> & \centering 2,0...50,0 & \centering В & \centering \arraybackslash 0,1 \\
\hline
\centering Контроль ТН (Контроль\_ТН)  & \centering SGF3 & \centering 3 фазы \\ 2 фазы & \centering -- & \centering \arraybackslash -- \\
\hline
\centering Режим пуска (Реж\_пуска)  & \centering SGF2 & \centering Измеренный \\ Внешний & \centering -- & \centering \arraybackslash -- \\
%===t1
\end{longtable}
\normalsize

\end{enumerate}
